\chapter{Explicit Algorithmic State in Uniaxial Steel and Concrete Laws}

\section{Why This Step Follows Naturally}

Chapter~31 proved that the constitutive-site state can be externalised for a
continuum inelastic law without breaking the legacy API.  The next
architecturally meaningful target is the pair of uniaxial nonlinear laws that
support reinforced-concrete fiber sections:
\texttt{MenegottoPintoSteel} and \texttt{KentParkConcrete}.

That propagation matters for two reasons:

\begin{itemize}
  \item reinforced concrete in \texttt{fall\_n} is already assembled from many
        uniaxial constitutive sites embedded in a section reduction;
  \item once those sites own their irreversible variables explicitly, the same
        architectural pattern scales from continuum plasticity to structural
        reduction and eventually to multi-scale transfer.
\end{itemize}

\section{Semantic Improvement}

Before this step, both uniaxial laws still mixed two different roles inside a
single object:

\begin{enumerate}
  \item the constitutive map itself, that is, the stress--strain law;
  \item the site memory needed to follow hysteresis, unloading branches,
        cracking, yielding, reversal points, and so on.
\end{enumerate}

That coupling was acceptable for a local prototype, but it was semantically
inconsistent with the direction already established for
\texttt{PlasticityRelation}.  A constitutive law should define the local map,
while the constitutive site should own the algorithmic state.

\section{New Contract}

Both \texttt{MenegottoPintoSteel} and \texttt{KentParkConcrete} now satisfy
\texttt{ExternallyStateDrivenConstitutiveRelation}.  Concretely, they support
the explicit-state interface

\begin{lstlisting}[language=C++]
ConjugateT compute_response(const KinematicT& eps,
                            const InternalVariablesT& alpha) const;

TangentT tangent(const KinematicT& eps,
                 const InternalVariablesT& alpha) const;

void commit(InternalVariablesT& alpha,
            const KinematicT& eps) const;
\end{lstlisting}

This means that the constitutive law can now be used as the pure local map
\[
(\varepsilon_{n+1}, \alpha_n)
\mapsto
(\sigma_{n+1}, C_{n+1}^{\text{alg}}, \alpha_{n+1})
\]
without requiring \(\alpha_n\) to live inside the relation object itself.

\section{Menegotto--Pinto Steel}

For \texttt{MenegottoPintoSteel}, the explicit state object stores:

\begin{itemize}
  \item the last reversal point;
  \item the current asymptote-intersection point;
  \item accumulated plastic excursion;
  \item loading direction and committed strain/stress.
\end{itemize}

The key refactor was to isolate the state update into a dedicated
\texttt{commit\_state()} kernel.  The explicit-state path now:

\begin{enumerate}
  \item evaluates the current branch using the supplied
        \texttt{MenegottoPintoState};
  \item commits the new reversal/branch information into that explicit state;
  \item leaves the embedded compatibility state untouched.
\end{enumerate}

The old API still exists:
\texttt{compute\_response(eps)}, \texttt{tangent(eps)}, \texttt{update(eps)},
and \texttt{internal\_state()} remain available as a compatibility layer.

\section{Kent--Park Concrete}

For \texttt{KentParkConcrete}, the explicit state object stores:

\begin{itemize}
  \item the most compressive strain reached on the envelope;
  \item the stress at that point;
  \item the unloading plastic strain;
  \item cracking status in tension;
  \item the last committed pair \((\varepsilon,\sigma)\).
\end{itemize}

Here the refactor isolated the old \texttt{update()} behaviour into a
\texttt{commit\_state()} function that operates on an explicit
\texttt{KentParkState}.  This makes the simplified compression-envelope /
unloading logic reusable by both:

\begin{itemize}
  \item the legacy embedded-state path;
  \item the new explicit-state path used by constitutive sites.
\end{itemize}

\section{Why \texttt{FiberSection} Did Not Need a New API}

This point is important architecturally.  \texttt{FiberSection} already stores
each fiber as

\begin{center}
\texttt{Material<UniaxialMaterial>}
\end{center}

which in turn delegates to a typed constitutive site
\texttt{MaterialInstance<Relation, StatePolicy>}.  Since
\texttt{MaterialInstance} now prefers the explicit-state path whenever the
relation satisfies \texttt{ExternallyStateDrivenConstitutiveRelation}, the
benefit propagates automatically to fibers.

Therefore no new visualization path, no new section API, and no new type
erasure was needed.  The existing architecture was already good enough; it
only needed the constitutive laws to satisfy the stronger contract.

\section{What the New Tests Prove}

The new uniaxial regressions verify the following facts:

\begin{enumerate}
  \item an external \texttt{MenegottoPintoState} evolves when committed
        explicitly, while the naked relation keeps its compatibility state
        untouched;
  \item an external \texttt{KentParkState} evolves when committed explicitly,
        and unloading from that state differs from the virgin-envelope
        response;
  \item a typed constitutive site
        \texttt{MaterialInstance<MenegottoPintoSteel, CommittedState>} stores
        the steel hysteretic state at the site level;
  \item a typed constitutive site
        \texttt{MaterialInstance<KentParkConcrete, CommittedState>} stores the
        concrete envelope/unloading state at the site level;
  \item existing fiber-section tests still pass, confirming that the new
        constitutive-site ownership model does not break the section reduction.
\end{enumerate}

\section{Architectural Consequences}

After this step the semantic picture is more coherent:

\begin{itemize}
  \item \texttt{MenegottoPintoSteel} and \texttt{KentParkConcrete} are
        constitutive laws with an optional embedded compatibility state;
  \item \texttt{MaterialInstance} is the real constitutive site;
  \item \texttt{EmbeddedRelationIntegrator} can now drive steel, concrete, and
        plasticity through the same explicit-state protocol;
  \item \texttt{FiberSection} remains a reduction policy over many constitutive
        sites instead of becoming yet another owner of hidden irreversible
        memory.
\end{itemize}

\section{What Still Remains to Be Done}

This chapter does \emph{not} complete the materials refactor.  The following
steps are still pending:

\begin{itemize}
  \item move the same pattern into section-level inelastic reductions where it
        is still mixed with law logic;
  \item let specialized constitutive integrators replace the currently embedded
        update rules when substepping or more elaborate algorithms are needed;
  \item introduce regularized damage / fracture models for continuum concrete
        so that cracking can be represented objectively with respect to the
        mesh.
\end{itemize}

\section{Strategic Importance for Reinforced Concrete Across Scales}

This propagation is small in terms of code size, but large in architectural
value.  Reinforced concrete across scales requires:

\begin{itemize}
  \item uniaxial steel and concrete at the fiber level;
  \item section reductions for beams, columns, and shells;
  \item continuum models with explicit kinematics, inertia, and possibly
        damage or fracture;
  \item future transfer of reconstructed fields across scales.
\end{itemize}

Externalising the algorithmic state at the uniaxial level is the first place
where these scales start to line up under the same constitutive-site
abstraction.  That is the real reason this step matters.
